\documentclass[11pt]{article}
\usepackage[margin=0.85in]{geometry}
\usepackage{amsmath,amssymb,mathtools}
\usepackage{booktabs,array,enumitem,microtype}
\setlength{\parindent}{0pt}
\setlength{\parskip}{0.65em}

\newcommand{\ket}[1]{\lvert #1\rangle}
\newcommand{\CNOT}[2]{\operatorname{CNOT}_{#1\to #2}}

\begin{document}

\begin{center}
  {\Large\bfseries Problem 1 -- Question 4}\\[0.3em]
  {\large State preparation}
\end{center}

\section*{4(a): Uniform superposition}

Apply a Hadamard gate to every qubit:
\[
 \ket{0}^{\otimes n}
 \xrightarrow{\,H^{\otimes n}\,}
 \left(\frac{\ket0+\ket1}{\sqrt2}\right)^{\otimes n}
 =\frac{1}{\sqrt{2^n}}\sum_{x\in\{0,1\}^n}\ket{x}.
\]
The circuit has size \(n\), depth \(1\), and uses no ancillary qubits.
Both the size and depth are optimal: every output qubit must be changed
from the fixed state \(\ket0\), and all \(n\) one-qubit gates can run in
parallel.

\section*{4(b): Unequal-amplitude GHZ state}

Assume \(|\alpha|\leq1\), and choose
\[
 \gamma=2\operatorname{atan2}\left(\sqrt{1-\alpha^2},\alpha\right).
\]
Then
\[
 R_Y(\gamma)\ket0
 =\alpha\ket0+\sqrt{1-\alpha^2}\ket1.
\]
Apply \(R_Y(\gamma)\) to \(q_1\), followed by a fanout of CNOTs whose
control information originates at \(q_1\):
\[
\left(\alpha\ket0+\sqrt{1-\alpha^2}\ket1\right)\ket{0}^{\otimes(n-1)}
\longmapsto
\alpha\ket{0^n}+\sqrt{1-\alpha^2}\ket{1^n}.
\]
With unrestricted connectivity, use a balanced doubling tree: in layer
\(r\), every already-entangled qubit controls a fresh qubit. This requires
\(n-1\) CNOTs and depth \(\lceil\log_2 n\rceil\), after the initial
one-qubit rotation. Hence
\[
 \text{size}=n,\qquad
 \text{depth}=1+\lceil\log_2 n\rceil,\qquad
 \text{ancillas}=0.
\]
On a nearest-neighbor line, the simple CNOT chain has the same size and
depth \(n\).

\section*{4(c): Four-qubit one-excitation state}

Use qubit order \(q_1q_2q_3q_4\). Define an excitation-preserving Givens
rotation by
\[
\begin{aligned}
G_{jk}(\varphi)\ket{01}_{jk}
 &=\sin\varphi\,\ket{10}_{jk}
  +\cos\varphi\,\ket{01}_{jk},\\
G_{jk}(\varphi)\ket{10}_{jk}
 &=\cos\varphi\,\ket{10}_{jk}
  -\sin\varphi\,\ket{01}_{jk},
\end{aligned}
\]
and let it act as the identity on \(\ket{00}\) and \(\ket{11}\).
Starting from \(\ket{0000}\), use the chronological sequence
\[
 X_{q_4},\quad
 G_{34}(\varphi_{34}),\quad
 G_{23}(\varphi_{23}),\quad
 G_{12}(\varphi_{12}),
\]
where
\[
 \cos\varphi_{34}=\frac{3}{10},\qquad
 \cos\varphi_{23}=\frac{4}{\sqrt{91}},\qquad
 \varphi_{12}=\frac{\pi}{4}.
\]
The amplitudes evolve as
\begin{align*}
\ket{0001}
&\longmapsto
\frac{\sqrt{91}}{10}\ket{0010}
+\frac{3}{10}\ket{0001}\\
&\longmapsto
\frac{\sqrt3}{2}\ket{0100}
+\frac25\ket{0010}
+\frac{3}{10}\ket{0001}\\
&\longmapsto
\frac{\sqrt6}{4}\ket{1000}
+\frac{\sqrt6}{4}\ket{0100}
+\frac25\ket{0010}
+\frac{3}{10}\ket{0001}.
\end{align*}
This is the requested state. It uses one \(X\) gate and three two-qubit
Givens rotations, no ancillas, and depth \(4\). A Givens rotation can be
compiled into two CNOTs and single-qubit \(R_Y\) pulses if CNOT is the
native entangling gate.

\section*{4(d): Optional part retained as requested}

Let the eight strings multiplying \(\alpha\) be the affine code
\[
 \mathcal A
 =a_0+\operatorname{span}_{\mathbb F_2}\{g_1,g_2,g_3\},
\]
where
\[
\begin{aligned}
a_0&=1000000,\\
g_1&=1010101,\qquad
g_2=0110011,\qquad
g_3=0001111.
\end{aligned}
\]
The eight strings multiplying \(\beta\) are exactly the bitwise
complements of the strings in \(\mathcal A\). Thus they form
\(\mathcal A+1111111\). The state is normalized when
\[
 8(\alpha^2+\beta^2)=1.
\]

Choose \(\delta\) such that
\[
 \cos\frac{\delta}{2}=\sqrt8\,\alpha,\qquad
 \sin\frac{\delta}{2}=\sqrt8\,\beta.
\]
Initialize seven qubits and prepare four logical bits
\[
 u,v,w\ \text{uniform in }\{0,1\},
 \qquad
 s\ \text{with amplitudes }\sqrt8\alpha,\sqrt8\beta,
\]
by applying
\[
 H_{q_1},\quad H_{q_2},\quad R_Y(\delta)_{q_3},\quad H_{q_4}.
\]
Initially \(q_1=u,q_2=v,q_3=s,q_4=w\), while \(q_5=q_6=q_7=0\).

Apply the following compact linear encoder in chronological order:
\[
\begin{aligned}
&\CNOT{1}{2},\
\CNOT{1}{5},\
\CNOT{1}{6},\
\CNOT{2}{7},\
\CNOT{3}{1},\
\CNOT{1}{2},\\
&\CNOT{3}{4},\
\CNOT{7}{3},\
\CNOT{4}{7},\
\CNOT{7}{6},\
\CNOT{4}{5}.
\end{aligned}
\]
Tracking the logical bits over \(\mathbb F_2\), the seven registers become
\[
\begin{aligned}
q_1&=u+s,&q_2&=v+s,&q_3&=u+v+s,\\
q_4&=w+s,&q_5&=u+w+s,&q_6&=v+w+s,\\
q_7&=u+v+w+s.&&&
\end{aligned}
\]
Finally apply \(X_{q_1}\). The output bit string is therefore
\[
 a_0+ug_1+vg_2+wg_3+s(1111111).
\]
Since \(H^{\otimes3}\) supplies the factor \(1/\sqrt8\), every word with
\(s=0\) has amplitude \(\alpha\), and every word with \(s=1\) has
amplitude \(\beta\), exactly as required.

This explicit encoder uses no extra ancillas. Counting the displayed
primitive gates gives \(3\) Hadamards, \(1\) \(R_Y\), \(11\) CNOTs, and
\(1\) \(X\), for total size \(16\). Disjoint CNOTs may be parallelized.

\end{document}
